% ============================================================
%  Harald Høffding, "Totalitet som Kategori:
%  en erkendelsesteoretisk Undersøgelse"
%  "Totality as a Category: An Epistemological Investigation"
%  Det Kgl. Danske Videnskabernes Selskabs Skrifter,
%  7. Række, historisk og filosofisk Afd. III, 2.
%  København: Bianco Lunos Bogtrykkeri, 1917.
%
%  TRANSLATION (English).  TRANSLATION-ONLY PROJECT: no Danish
%  transcription is produced or published.
%
%  ---------------------------------------------------------------
%  SOURCES AND THE PAGE MAP
%  ---------------------------------------------------------------
%  Base text: the DRM-free Danish e-book (SAGA / Lindhardt og
%  Ringhof, ISBN 9788726364453), which is an OCR of the 1917 print.
%
%  Page-accurate witness: the Royal Danish Academy's digitisation of
%  the 1917 printing, <http://publ.royalacademy.dk/books/419/2912>,
%  on disk as ~/bibliotek/Høffding, Harald/totalitet.pdf (80 pp.,
%  ABBYY text layer).  THE TREATISE'S OWN PRINTED PAGE NUMBER EQUALS
%  THE PDF PAGE NUMBER (offset 0).  Every page also carries the
%  containing volume's continuous pagination, which is the printed
%  page + 256; the treatise occupies volume pp. 259--334.  Høffding
%  cites this treatise by its own pagination (3--78) in the 1925
%  "Erkendelsesteori og Livsopfattelse", so 3--78 is the numbering
%  used in the markers below.
%
%  The e-book's endnotes carry NO printed page references, unlike
%  those of "Filosofiske Problemer"; the page map here therefore
%  comes from the print itself, not from a reconstruction.  It is
%  confirmed twice over: by the running heads, and by the book's own
%  INDHOLD on p. 80.
%
%  ---------------------------------------------------------------
%  STRUCTURE
%  ---------------------------------------------------------------
%  Six chapters, and TWENTY-NINE numbered run-in divisions that run
%  STRAIGHT THROUGH the whole treatise (1--29) rather than restarting
%  at each chapter.  They are set here with \hnum, as in the
%  translation of the 1925 "Erkendelsesteori og Livsopfattelse".
%
%    ch.  pp.     divisions (first printed page of each)
%    I    3--23   1 p3, 2 p4, 3 p6, 4 p8, 5 p13, 6 p17, 7 p18, 8 p20
%    II   23--31  9 p23, 10 p28
%    III  31--42  11 p31, 12 p32, 13 p34, 14 p38, 15 p39
%    IV   42--66  16 p42, 17 p45, 18 p48, 19 p50, 20 p52, 21 p56,
%                 22 p59, 23 p63
%    V    66--73  24 p66, 25 p68, 26 p70, 27 p72
%    VI   74--78  28 p74, 29 p75
%    Appendix (KATEGORITAVLE)  p. 79
%
%  ---------------------------------------------------------------
%  THE CONTENTS PAGE DISAGREES WITH THE CHAPTER HEADS
%  ---------------------------------------------------------------
%  Recorded, not normalised, per the repo's editorial stance.  The
%  chapter heads (pp. 23, 31, 42, 66) read "Totalitetsbegrebet og de
%  fundamentale / formale / reale Kategorier" and "Totalitet og
%  Værdi"; the INDHOLD on p. 80 reads "Totalitet og de fundamentale
%  / formale / reale Kategorier" and "Totalitet og Værdikategorierne".
%  The chapter heads are followed here, as the fuller reading and the
%  one at the head of the text itself.
%
%  ---------------------------------------------------------------
%  THE PRINT'S OWN ERRATA LIST (TRYKFEJL, p. 80)
%  ---------------------------------------------------------------
%  Eight corrections, all in chapters III--IV.  Each is to be applied
%  when its page is reached, and logged in RESUME-NOTES as applied:
%    p. 31, l. 6 from foot:  deres            -> den
%    p. 33, l. 21 from top:  Individualforestillinger -> Individualbegreber
%    p. 36, l. 1 from top:   fra              -> for
%    p. 36, l. 20 from top:  Erfaringer       -> Erfaringen
%    p. 36, 2nd note, l. 3:  samtlige         -> sanselige
%    p. 44, l. 19 from top:  Farveaktens      -> Farvesansens
%    p. 45, 1st note:        340              -> 388
%    p. 64, l. 3 from foot:  første           -> faste
%  None falls in the pages translated so far.
% ============================================================

\documentclass[12pt]{book}
\usepackage[utf8]{inputenc}
\usepackage[T1]{fontenc}
\usepackage[english]{babel}
\usepackage[protrusion=true,expansion=true]{microtype}
\usepackage[margin=1.3in]{geometry}
\usepackage{setspace}
\usepackage{amsmath}
\usepackage{enumitem}
\usepackage{libertinus}
\usepackage{libertinust1math}
\usepackage{textalpha}
\usepackage{fancyhdr}
\usepackage[colorlinks=true,linkcolor=black,urlcolor=black,
  pdftitle={Totality as a Category},
  pdfauthor={Harald Høffding}]{hyperref}
\renewcommand{\thechapter}{\Roman{chapter}}
\setlength{\headheight}{15pt}
\pagestyle{fancy}
\fancyhf{}
\fancyhead[LE,RO]{\thepage}
\fancyhead[RE]{\textit{Totality as a Category}}
\fancyhead[LO]{\textit{\leftmark}}
\setlength{\parindent}{1.5em}
\setstretch{1.3}

% run-in division number (Høffding's numbering, continuous 1--29)
\newcommand{\hnum}[1]{\textbf{#1.}\enspace}

\title{Totality as a Category\\[0.5em]
  \large An Epistemological Investigation}
\author{Harald Høffding}
\date{Translated from the Danish\\[0.6em]
  \small \textit{Det Kgl.\ Danske Videnskabernes Selskabs Skrifter},
  7.\ Række,\\
  \textit{historisk og filosofisk Afd.} III, 2.\\
  Copenhagen: Bianco Lunos Bogtrykkeri, 1917}

\begin{document}
\maketitle
\tableofcontents


% ============================================================
%  I. THE POSITION OF THE DOCTRINE OF CATEGORIES IN PHILOSOPHY
%     pp. 3--23; divisions 1--8; notes 1--12
% ============================================================

\chapter{The Position of the Doctrine of Categories in Philosophy}
\label{ch:1}

\hnum{1} Categories are fundamental concepts which are contained in the
presuppositions on which all scientific research rests, and which can be
exhibited by analysis of those presuppositions. But for a full elucidation of
their nature it is not enough to start out from the science existing at a given
time. It must also be shown how they have worked their way forward through the
history of science, --- how they arise as forms of thought-work, --- and how
they can disappear again when they are no longer expressions of labouring
thought. And from such a historical consideration one must go back again to a
psychological investigation of the development of thinking in ordinary human
conscious life, of which scientific thinking is a peculiar formation conditioned
by determinate tasks. The fundamental concepts with which science works must be
psychologically possible, and their first origin must be findable in involuntary
thought-life. The difference between the psychological and the epistemological
consideration of the categories consists in this, that the former finds
possibility for unscientific as well as for scientific thinking, whereas the
latter shows that under certain given conditions there is only one single narrow
road to understanding. Of our thoughts it holds that many are called, few
chosen, --- indeed, we must add as psychologists, that many are uncalled.

In epistemological works one will find different weight laid on fundamental
concepts and on fundamental propositions. Some epistemologists keep solely to
the fundamental propositions, the propositions on which scientific grounding
rests. And since propositions that are presuppositions for all grounding
naturally cannot themselves be grounded, the fundamental propositions are often
presented as resting on a choice, a leap, an act which in and for itself might
just as well have fallen out in quite another direction. The first principles of
science are, on this conception, not themselves an object for science. Yet it is
for the most part conceded by those who hold this conception that the choice is
bound within a certain circle of possibilities, and that this circle of
possibilities is given with the nature of human thinking, such as this once and
for all is constituted. It is psychology's task to investigate this factual
constitution, and the question then presents itself how the different
standpoints from which epistemology and psychology assess human cognition stand
to one another. But even before we take up this question it will be clear that
there is such a reciprocal relation between fundamental concepts and fundamental
propositions, as in general between concepts and judgments, that the
investigation of the single concepts which the fundamental propositions contain
will throw light on the fundamental propositions themselves. Advancing clarity
becomes possible at all only by our investigating and determining, with express
consciousness, the concepts contained in our judgments, just as on the other
side we investigate and determine the kind of connection that is asserted in the
judgments between the concepts. The two investigations supplement one another.
Advancing clarity in the determination of concepts is the condition for
advancing clarity in the understanding of the judgments into which the concepts
enter, and conversely. Whether one lays the greatest weight on the one or the
other investigation will depend on the special task one sets oneself. But for a
complete reckoning neither of the two standpoints can be dispensed with. To this
it comes in addition that the type of thought is the same in concepts as in
judgments. The judgment is a connection of concepts, but the concept is a
connection of marks. It is the same fundamental property of our thinking that
expresses itself in both places, in the concept more as structure, in the
judgment more as function.

Whereas by keeping merely to science's first presuppositions one arrives at an
uncompromising theory of postulates, which places all scientific thought-work in
a distant relation, or even in sharp opposition, to all other spiritual work,
one will, by investigating the fundamental concepts as well, be able to be led
to an understanding of the relation of scientific thought-work to the life of
the spirit as a whole. And since it is precisely one of philosophy's tasks to
investigate this relation, it will be of importance to attempt a determination
of the position of the doctrine of categories within philosophy.

\hnum{2} When does science begin? And how does scientific consciousness stand to
the thought-life that unfolds before and alongside the scientific, in greater or
less independence of it?

For the elucidation of these questions it must first be held fast that a common
type makes itself felt in all human thought-life, the prescientific, the
scientific and the unscientific alike.

Practical life itself, its necessities and its tasks, leads to thinking. In
order that a man may attain the ends he strives after, there are certain means
that must be procured. Here man learns in fact and in practice what necessity
is, that is, he learns to know the inexorable relation between ground and
consequent. So long as he holds fast to the end, so long must he also concede
the necessity of means, and finally of certain determinate means; and he learns
conversely to see that the attainment of the end is a consequence of the
application of certain means. He learns besides that when the same necessities
and the same outer conditions repeat themselves, the same means must be resorted
to. A Central Australian, for example, gathers plants, threshes them with a
stick, sifts the seed, roasts it, crushes it on stones and eats it: a series of
actions which, every time he is hungry, must be performed in the same way and in
the same order. At bottom, science's first presuppositions may be said to be
present already in this series-formation.

In the most recent time particular weight has been laid on the influence of
social life on the first development of thinking. In the tribe's division into
different groups or classes there lay a model for all classification, and
thereby for all logic. We still speak, after all, of genera and families in the
division of natural objects, in spite of the fact that eager theorists of
heredity find false analogies here. The representation of time and the division
of time received their development through the rhythmical course and regular
recurrence of religious feasts and ceremonies. Our almanac still indicates the
scheme of the church year. And religious worship demanded that liturgical acts
be performed in a quite determinate manner, the same in every single case, if
the result was to be attained. Here, then, was inculcated the proposition that
like must be produced by like. The division of space received its clarity by
reason of the religious significance of the different directions and the
different places.\footnote{\textsc{Émile Durkheim} has, especially in his work
\textit{Les formes élémentaires de la vie religieuse} (1912), made this way of
looking at the matter prevail and sharpened it into the assertion that all
categories are of religious-social origin. Cf.\ my treatise \textit{Sociologi og
Religionsfilosofi} (Vor Tid, 1914--15), (in French in \textit{Revue de la
Métaphysique et de Morale}, Nov.\ 1914).}

That the social conditions under which men develop exercise a very great
influence on their way of thinking and their forms of thought is beyond doubt,
even if not everything can be explained along this road. Social influence is
constantly broken against the practical experience individuals can make on their
own account, face to face with those conditions of life which also gain
influence on the conduct of society as a whole.

Whether one will now lay the greatest weight on individuals' practical
experiences or on the influence of social institutions, it is a fact that at
each given time and among each single people there appears a certain conception
of the world, involuntarily and in part unconsciously developed, which is built
after the same main forms as those that appear in a more stringent and special
manner in scientific thinking. It bears witness to a need and a capacity to
order and interpret observations and traditions and thereby to bring them into a
certain connection. The thought-work that is done here proceeds without clear
consciousness. No attention is directed to the single transitions of thought or
to the warrant for connecting, in a certain determinate way, the single elements
that observation and tradition offer. To use the terminology I have employed
elsewhere: it is not analytic and synthetic intuition, but practical and
concrete intuition, that is present here.\footnote{See \textit{Om Begrebet
Intuition} (Vid.\ Selsk.\ Oversigt 1914), p.\ 77.} So-called sound common sense
operates throughout with these last-named kinds of intuition. Through them it
can, to use \textsc{Meyerson}'s expression,\footnote{\textsc{Meyerson} has
discussed the concept of \textit{sens commun} in an excellent manner in his book
\textit{Identité et Réalité} (1908), Chap.\ XI.} give ``a first, very clumsy
sketch of a scientific and metaphysical system''. --- In language, in the words
language forms and the word-combinations it shapes, such a system is already
contained, one that constantly makes its influence felt on the scientific forms
of thought.

The origin and development of science, as the history of science shows it to us,
really makes the question of its first beginning fall away entirely, since it
turns out that its first germ is present as soon as men begin to think. At all
stages and under all conditions a certain opposition makes itself felt between
something given, something presented, and its interpretation, between the single
subject-matter and the connection into which it is fitted. Progress consists in there
coming to be an ever closer and stricter relation between observation and
interpretation. More and more one becomes conscious that all interpretation must
in the end itself be fetched from observations, and that the fitting of the
single subject-matter as a member into a whole connection must take place after exact
testing both of the single observation itself and of the other observations with
which it is connected. As will be shown in what follows, it is here especially
the difference between identity and analogy, and the limit of analogy's warrant,
that one must have an eye for.

\hnum{3} It is psychology's task to find the common type for all thought-life,
the scientific as well as the prescientific and the unscientific.

This type can in its generality be designated as a whole-formation. Even the
simplest sensations of sense show themselves to stand in connection with one
another, so that the single one cannot be understood, either as regards degree
of strength or as regards quality, except by being seen as a member in a whole
series or group of sensations. The temporal relation between the different
sensations is already decisive for the degree of independence with which they
appear. To this it comes in addition that reproductions of earlier sensations
unite themselves more or less closely with the newly emerging sensations, often
from the very hour of their birth. Thereby recognition, partial perception and
sense-illusion become possible. Where the reproduced elements gain the decided
upper hand, memory and imagination appear as a whole-formation of a higher
order. Through association, displacement and comparison, intuitions of sense, of
memory and of imagination can undergo alterations, articulations, without the
character of intuitability disappearing. Cooperating in all these
whole-formations, from first to last, are practical motives, need and impulse,
pleasure and displeasure.

Characteristic of the whole-formations here spoken of (which lead to what I have
called concrete and practical intuition) is the involuntary manner in which they
come to be. From this follows the self-evidence and originality they seem to
possess. The elements of which they are built up are not themselves accessible
to us in the same way as the elements out of which we form concepts, judgments
and inferences with full consciousness, since the marks of the concept, the
members of the judgment and the premises of the inference are then given in
advance of the new whole-formations into which they enter. Thereby what I have
called analytic and synthetic intuition stands in opposition to concrete and
practical intuition. Now it is in fact not all concepts, judgments and
inferences that we form with full consciousness; they can also arise by
involuntary whole-formation. But if concepts, judgments and inferences are to
have scientific significance, they must be capable of being formed anew, in such
a way that full and clear consciousness is present at each single transition
from thought to thought (analytic intuition) as well as at each single
connection of thought with thought (synthetic intuition). Yet even in the
involuntary formation of concepts, judgments and inferences, reflection appears
as something different from immediate intuiting (in sensation, memory and
imagination), inasmuch as a more or less distinct transition and comparison
takes place between the elements that are brought into connection, whereas the
elements out of which immediate intuitions are formed can be traced only
indirectly and in part by the road of analogy.

Since the formation of concepts takes place by judging, and since the inference
is the formation of a judgment out of other judgments, we can in this connection
keep to the judgments.

The judgment is formed in the first instance on the basis of immediate
intuiting. It states a determinate relation between elements in the content of
the intuition, elements which in the intuition itself stand in continuous
connection or are perhaps even fused with one another. The judgment ``the table
is square'' is formed on the basis of an intuiting in which the property
``square'' was already contained (``lay''), but without being specially brought
forward and without being placed in relation to the table as a whole, that is,
as the sum-total of all its properties.

But why are judgments formed at all, --- why do we not remain at intuiting? It
is surely not, as some authors seem to conceive the matter, a kind of fall into
sin that has taken place here. It is the richness and manifoldness of
observation, and memory's attempt to span this richness, that make another kind
of whole-formation than immediate intuiting necessary. The primary form of the
whole is burst open, and a new one is sought.

It can be difficult to draw a sharp boundary between intuition and judgment. As
already indicated, the intuition can be altered, articulated. This happens by
single parts pushing themselves into the foreground while others recede. In the
case of the table, for example, the colour can make itself specially felt in
preference to the figure, and in the figure itself the contour in preference to
the form of the surface.\footnote{Cf.\ \textsc{Edgar Rubin}: \textit{Synsoplevede
Figurer}.} Such displacements go on unceasingly in our sensations, memories and
imaginings. But a judgment is formed only when we become conscious of these
transitions instead of being wholly absorbed in the intuitions that are present
in the single moments and that succeed one another. Now as regards this
consciousness there may indeed be many transitions and degrees; but the
end-points nevertheless stand in characteristic opposition to one another, and
it would be unwarranted to transfer to intuition everything that holds of
judgment. There is only analogy between them, an analogy grounded in their being
two different ways in which psychical energy, that is, the need and the capacity
for whole-formation, can express itself.

The founders of the doctrine of categories, \textsc{Aristotle} and
\textsc{Kant}, rightly saw that the scientific fundamental forms of cognition
must be findable by analysis of the judgments that are passed. What they have
not seen clearly enough is that the psychical energy at work in the formation of
judgment also works, under other conditions, in other whole-formations. Only
through the continuity and analogy that make themselves felt between the primary
whole-formation and the judgment is it understood how scientific consciousness
stands in close connection with the life of consciousness as a whole and is a
more special formation of it. But it is also precisely only a relation of
analogy that is present, as between two examples of the same concept. Since
judgments can be derived from, built upon, intuitions, one may readily say that
they are ``contained'' in them or ``lie'' in them. But that is figurative
speech. For because a judgment is logically possible, it is not thereby said
that it is psychologically possible; and in any case it is a conversion into
another form that takes place in the transition from intuition to judgment, a
transition that presupposes certain determinate conditions. One must not confuse
the immediate and involuntary security with which the mind rests in intuitions
of sense and of memory with the assurance that is first the work of reflection.
Such a confusion is committed by \textsc{Heinrich
Maier}\footnote{\textit{Psychologie des emotionalen Denkens} (1908), p.\ 2. ---
Cf.\ on this whole question \textit{Den menneskelige Tanke}, p.\ 53--59; 74 f.;
99--103.} when he declares: ``Elementary judgments lie in every observation,
every memory-image, every cognitive representation of imagination, in so far as
all these representations have real objects, real events, states, things and so
forth for their object''. And when he adds: ``The alleged images of observation,
memory and imagination, which are supposed to be indifferent to the opposition
between true and false, valid and invalid, are fictions that never and nowhere
are real'', --- then I can only agree with him in so far as that immediate
security cannot be called indifference. Strictly speaking one can speak of
indifference only where both members of the opposition are known and stand as
equally valid. But this acquaintance is not present from the first. There has
not yet been any eating of the tree of knowledge. No doubt has yet stirred; the
state is positive through and through. When the immediate security is broken
through contradiction, or through the demands for closer determination that new
observations make, then reflection begins its work of producing new
whole-formations, and then the judgment can replace the intuition. Only then
does application arise for the concepts true and false as opposites.

\hnum{4} The difference between elementary, involuntary thought-work and
scientific thought-work can be elucidated from three different sides.

\textit{a.} Reflection will first of all seek to order the different
subject-matters it distinguishes according to their mutual relation. Relation,
determinateness by relation, is therefore (as the history of the doctrine of
categories also shows) the most conspicuous category. But orderings determined
by relation can be of different significance for the development of cognition.
Through a series of stages reflection advances from purely external orderings,
which cannot lead further, to orderings that make it possible to derive new
relations from those already present. There is a whole series of stages from
chaos to rationality, and rationality can itself have different degrees. These
stages are elucidated by the series that can be formed by the ordering of the
subject-matters. A survey of these series I have given in ``Den menneskelige
Tanke'' (p.\ 162--207); but I come back to them here in all brevity and from a
somewhat different point of view.\footnote{As an appendix to the present
treatise there is printed a list of the categories, essentially as I have
presented them in the book named. And there are also set out the different
series-formations that express the transition from prescientific to scientific
consciousness.}

α. If we try to think for ourselves the conditions most unfavourable for human
thinking, we stand before the logically irrational. We stand before such a
relation between given subject-matters that they are all equally different from
one another, and that they can therefore be placed in any order whatever. One
may push out any member whatever without the relation between the remaining
members being altered, and without anything following from it at all. Such a
series I have called a chaotic difference-series. It does not appear in
immediate observation. For already, purely psychologically, the order of the
addends is not a matter of indifference. Our mental states are determined, both
as regards degree and as regards quality, by their place in the series of our
lived experiences. But there can be all possible approximations to such a
logical irrationality, which differs from mathematical irrationality in this,
that it cannot, as the latter can by its very manner of coming to be, lead to a
reduction of the differences. We know that $\sqrt{2}$ lies in the series of
numbers between 1 and 2 and nowhere else, and within the interval named the
place can be determined ever more exactly the more decimals we calculate. The
logically irrational, by contrast, means a halt.

The next step is that series are formed whose members lie fast, while the
relations of difference vary from member to member. Such an indeterminately
varying difference-series does not permit elimination without the members
between which the elimination takes place coming into another relation to one
another than before. But the advance consists in this, that clear survey has
replaced the chaotic unrest.

A new advance is made when the differences between the members vary in a
determinate direction within the series. Such a regularly varying
difference-series makes possible for reflection the formation of a law or rule
by which the place of the subject-matters in the series can be determined. And
the task can be set of finding new members in the series, intermediate members
between those already given.

If the difference between members that have their determinate place in a series
is constantly the same, member for member, we get an identically varying
difference-series. The concept of identity begins to dawn, provisionally as
identity of the difference between qualitatively varying subject-matters. Such
series can be formed of moments, numbers, degrees and places. Every moment, with
its peculiar lived experience, has a determinate place in the series of moments,
and every moment stands in the same relation to the preceding and to the
following moment. So too with the numbers. Every number has from the first, as
\textsc{Poincaré} has expressed it, its own individuality, is qualitatively
different from other numbers, not merely by its place in the number-series. That
one can by different roads come to the same individual number is a discovery
that presupposes a higher stage than the one we here stand at. \textsc{Goethe}
held that mathematics moves in nothing but identities, because 2 + 2 = 4, but
also 4 = 2 + 2, so that we get nothing other than what we had beforehand. But
the different roads by which one can come to a number (e.g.\ by counting to 4 or
by adding 2 and 2 together) are different operations, of which one discovers
only afterwards that they lead to the same result and presuppose the same
principles. The number is provisionally an ordinal number, a numeral.
Differences of degree too are immediately differences of quality, and yet they
lie in a determinate order in a series whose members can vary in an identical
manner. So it stands with places in space as well. Even if they are perhaps
still regarded as qualitatively different, it is an identical varying by which
one passes from place to place.

β. In the series so far spoken of, the members are qualitatively different, and
it is merely a question of their mutual relation with respect to place in a
series. It is a great advance when this place becomes quite determinate.
Classification and comparison become possible thereby. But these operations can
really be carried out only when series can be formed in which every member is
comprised in the preceding one. This can again appear in two ways. In a
progressive difference-series every member can, for example, be a part of the
preceding. In a classification the species is a part of the genus, and the
individual a part of the species. From this it can be inferred that when an
individual \textit{A} belongs to the species \textit{B}, and the species
\textit{B} to the genus \textit{C}, \textit{A} also belongs to \textit{C}. It is
essentially on this that the Aristotelian doctrine of inference builds. In a
partial identity-series the following member stands to the preceding as
predicate to subject, or as consequent to ground. Here too classification can be
used as an example, for that \textit{A} belongs to the species \textit{B} means
that the predicates \textit{B} designates belong to \textit{A}. In both the
series named, the inference takes place by the elimination of intermediate
members.

When a series in which the members lie fast is symmetrical, that is, can be read
backwards just as well as forwards, one can infer from it to the reversed
series.

γ. Progressive difference-series, partial identity-series and symmetrical series
may be called rational series, inasmuch as they make inference possible. But the
qualitative difference of the members is not abolished by their being capable of
being seen as parts or wholes, or as subjects and predicates in relation to one
another. The most perfect kind of inference becomes possible only when not
merely the relation between the members is constantly identical, but the members
themselves are so too, so that not merely members can be pushed out, but every
member can be substituted for every other. We then get an absolute
identity-series: \textit{A} = \textit{B} = \textit{C} = \textit{D} \dots\ Here
is given, as \textsc{Leibniz} first rightly saw, the foundation --- and at the
same time the ideal --- for all rational thinking. Such a series designates the
complete opposite of the chaotic difference-series. The two series designate the
ideal limits between which human cognition moves.

Step by step we thus come, by running through the stages the different series
designate, up to the presuppositions on which strict science rests. But we see
at the same time that the rational series are not the only possible ones, and
that within the rational series a distinction must be made between those that
make merely classification and comparison possible and those that make inference
in the strictest sense possible. The different series form a scale, and it will
perhaps be possible to find still more stages on this scale than have been named
in the foregoing. The thought-work that is directed to the abolition of the
chaotic relation between the subject-matters presents a whole series of degrees,
until the purely rational stage is reached, where not merely the relation of the
members but the members themselves are identical. This highest stage stands,
however, always only as an ideal that can be reached only approximately and by
abstracting from a quantity of actually present differences. Through ordering,
classification and comparison, scientific consciousness works its way up towards
its highest form, everywhere that the subject-matters make it possible.

\textsc{Kant} took account, in ``Kritik der reinen Vernunft'', only of the
highest stage of science, that is, of the stage where it was possible to draw
inferences on the basis of absolute identity. Classifying and comparing science
he took no account of; that he reckoned to mere empiricism. And against the
objection that classification and comparison nevertheless presuppose the
validity of logical thinking no less than mathematical natural science does, he
answers that the principles of pure logic are indeed negative criteria of truth,
but not sufficient to ground the kind of science whose possibility he wishes to
investigate.\footnote{\textit{Kritik der reinen
Vernunft}\textsuperscript{1}, p.\ 151 f.} But we do after all ascribe objective
validity to our classifications and comparisons, just as much as to the natural
laws proper, and there is therefore just as much ground for investigating the
warrant for assuming the objective validity of purely logical principles as
there is for investigating the objective validity of the mathematical principles
and of the principle of causality. \textsc{Salomon Maimon} already saw this, in
maintaining against \textsc{Kant} that the logical fundamental forms,
affirmation and negation, have no significance if they are not presupposed valid
for the subject-matters to which they are applied. Whereas \textsc{Kant} lets
epistemology (``transcendental philosophy'') presuppose formal logic as a
discipline finished once and for all, \textsc{Maimon} maintains that epistemology
must precede formal logic if the latter is to have real
significance.\footnote{\textit{Versuch einer neuen Logik oder Theorie des
Denkens}, Berlin 1794, p.\ 404--408.} --- \textsc{Kant} naturally had the right
to limit the scope of his investigations, but he had not the right to think that
he had thereby also limited the scope of science.

\textit{b.} From a new side the transition from elementary, prescientific
consciousness to scientific consciousness comes forward when we bring out a
presupposition on which all the kind of series-formation spoken of in the
foregoing rests. We have in the foregoing taken account only of the struggle
against the differences, and have seen how the degrees of likeness, both as
regards relation and as regards members, step by step approach pure identity.
But all difference and likeness, and thereby all series-formation, presupposes
the possibility that subject-matters can be gathered, comprehended and ordered.
Here the concept of synthesis comes forward as a presupposition that holds both
for formal and for real science, both for classification and for the cognition
of laws. Series-formation rests on determinate relations between the
subject-matters being exhibited, and relation is thought only when the members
between which the relation holds are the object of one and the same act of
thought. Relation presupposes synthesis. But since synthesis,
gathering-together, consists precisely in bringing subject-matters into some
relation or other to one another, relation and synthesis are correlative
concepts.\footnote{In \textit{Den menneskelige Tanke} p.\ 155 f.\ I set up
synthesis as the first category, relation as the second. An acute reader has
drawn my attention to the fact that they ought to have been set up as
correlative concepts, like continuity and discontinuity, likeness and
difference.}

Now the series spoken of in the foregoing are only typical examples of what can
find application in many special cases. Many different classifications can be
undertaken, many different comparative points of view adopted, many different
laws found. But all these must again be brought into connection with one
another, and hereby new series-formations will come forward. Syntheses of a
higher order are then presupposed than in the single series taken each by
itself. What matters then is that the scope of the synthesis becomes as great as
possible. A striving towards an ideal maximum makes itself felt here. There will
here be greater resistance to overcome than in the single series, since the
difference between the series that can be formed in the different domains will
be greater than the difference between members in the same series and in the
same domain. The question becomes finally whether our thinking has at its
disposal logical forms for complete gathering-together, or whether at a certain
point more or less bold and more or less poetic analogies do not replace logical
thinking. In the sixth chapter of the present treatise we shall return to this
question. But before that boundary is reached where thinking passes over into
poetry, there is still a domain in which scientific hypotheses move with full
right. The need for gathering-together will constantly make itself felt, not
merely towards subject-matters, but also towards series. And by attempts at an
ever more comprehensive synthesis light can be thrown on the single series and
the single subject-matters, just as the single subject-matters are already
illuminated by the place they occupy in the series of which they are members.
Synthesis and analysis (and analogously: continuity and discontinuity, likeness
and difference) stand in constant reciprocal action with one another, and
science's rhythmical oscillations find their natural explanation in this
fundamental law of human thought.

\textit{c.} A third point of view for characterizing the difference between
prescientific and scientific thinking is given in this, that the more science
develops, the greater the differentiation of the different problems that takes
place. In involuntary thinking no distinction is made between subject-matter,
form and value. One does not make critically clear to oneself where the boundary
lies between the immediately given (the subject-matter) and the
thought-elaboration one applies in interpreting it. Nor is it recognized how
even the subject-matter's first appearance for us is conditioned by our
receptivity, which is never absolutely passive. Neither is value definitely
separated from subject-matter and form. The point of view of value will under
primitive conditions have been the predominant one, since subject-matters are
noticed especially on account of their practical significance, and since the way
in which the subject-matter is interpreted is also conditioned by its connection
with practical interests. \textsc{Bergson} ought therefore to have reckoned the
element of value among ``les données immédiates''; he has attempted to eliminate
the element of form and thinks that if it could be separated out, the element of
subject-matter alone would remain. He is even inclined to think that the ground
for the predominance of the element of form in science must be sought in
practical considerations. But interest cooperates already in the choice of
subject-matter. Naturally it cooperates also in the primitive interpretation.
The teleological point of view rules in animism under all its forms, from
fetishism right up into the most sublime theology. And in the primitive forms or
intimations of what we now call categories it can be traced that practical
motives have cooperated in their arising. Space means at first distance from an
end striven after, time a relation between a wished-for future and an
unsatisfying present, and the causal relation makes itself felt in the attempt
to abolish these distances in space and time. In the course of science's
development a constant striving is traceable to eliminate the element of value,
to abstract from all interest with the exception of the intellectual interest,
which both as motive and as effect is bound in the most intimate way to
thought-work itself. Thereby other values do not necessarily fall away, but the
concept of value can now become the object of a special problem, which among
other things also discusses the relation between intellectual value and other
values.

The distinction between the three elements marks an important advance. It has
not yet been strictly carried through in science, and still less in ordinary
consciousness. This distinction naturally solves no problems. On the contrary,
it is the condition for problems being able to be set sharply and determinately,
and the very clash between the three elements sets in the end the greatest
problem of all. --- For the rest, this point of view, drawn from
differentiation, leads back to the preceding points of view, inasmuch as with
increasing differentiation increasing demands are made on synthesis, if the
unity of the life of the spirit is to be preserved, and inasmuch as
differentiation occasions new relations.

% [text to be added: pp. 13--17 --- division 5, psychology's independence
%  over against epistemology; Kant's synthetic vs. analytic procedure.]

% [text to be added: pp. 17--20 --- divisions 6 and 7.  Division 6 ends with
%  the displayed series S₁{S₂{S₃{S₄ ……. (LEFT BRACES, as in "Filosofiske
%  Problemer" II §5 --- verified there on the 1902 page image; do not emend to
%  "<").  Division 7 has the displayed category-series
%  "Syntese > Identitet > Rationalitet > Kausalitet > Totalitet > Værdi",
%  where ">" is Høffding's own sign for an inconvertible, transitive and
%  asymmetrical relation (defined in "Den menneskelige Tanke" §66).]

% [text to be added: pp. 20--23 --- division 8, the historical career of the
%  concept of totality; notes 10--12; closes with four lines of Goethe quoted
%  in German (left in German, as printed), then a short centred rule.]


% ============================================================
%  II. THE CONCEPT OF TOTALITY AND THE FUNDAMENTAL CATEGORIES
%      pp. 23--31; divisions 9--10; notes 13--21
% ============================================================

\chapter{The Concept of Totality and the Fundamental Categories}
\label{ch:2}

% [text to be added: pp. 23--28 --- division 9; notes 13--17]

% [text to be added: pp. 28--31 --- division 10; notes 18--21.  The print's
%  errata list corrects p. 31, l. 6 from foot: "deres" -> "den".]


% ============================================================
%  III. THE CONCEPT OF TOTALITY AND THE FORMAL CATEGORIES
%       pp. 31--42; divisions 11--15; notes 22--45
% ============================================================

\chapter{The Concept of Totality and the Formal Categories}
\label{ch:3}

% [text to be added: pp. 31--34 --- divisions 11 and 12; notes 22--26.
%  Errata: p. 33, l. 21 from top, "Individualforestillinger" ->
%  "Individualbegreber".]

% [text to be added: pp. 34--38 --- division 13; notes 27--33.  Errata on
%  p. 36: l. 1 from top "fra" -> "for"; l. 20 from top "Erfaringer" ->
%  "Erfaringen"; 2nd note l. 3 "samtlige" -> "sanselige".]

% [text to be added: pp. 38--39 --- division 14; notes 34--38]

% [text to be added: pp. 39--42 --- division 15; notes 39--45]


% ============================================================
%  IV. THE CONCEPT OF TOTALITY AND THE REAL CATEGORIES
%      pp. 42--66; divisions 16--23; notes 46--84
%      The longest chapter: 69.8k characters, 5--6 batches.
% ============================================================

\chapter{The Concept of Totality and the Real Categories}
\label{ch:4}

% [text to be added: pp. 42--45 --- division 16; notes 46--50.  Errata:
%  p. 44, l. 19 from top, "Farveaktens" -> "Farvesansens"; p. 45, 1st note,
%  "340" -> "388".]

% [text to be added: pp. 45--48 --- division 17; notes 51--52]

% [text to be added: pp. 48--50 --- division 18; notes 53--58]

% [text to be added: pp. 50--52 --- division 19; note 59]

% [text to be added: pp. 52--56 --- division 20; notes 60--65]

% [text to be added: pp. 56--59 --- division 21; notes 66--70]

% [text to be added: pp. 59--63 --- division 22; notes 71--78]

% [text to be added: pp. 63--66 --- division 23; notes 79--84.  Errata:
%  p. 64, l. 3 from foot, "første" -> "faste".]


% ============================================================
%  V. TOTALITY AND VALUE
%     pp. 66--73; divisions 24--27; notes 85--97
% ============================================================

\chapter{Totality and Value}
\label{ch:5}

% [text to be added: pp. 66--68 --- division 24; notes 85--86]

% [text to be added: pp. 68--70 --- division 25; notes 87--91]

% [text to be added: pp. 70--72 --- division 26; notes 92--94]

% [text to be added: pp. 72--73 --- division 27; notes 95--97]


% ============================================================
%  VI. TOTALITY AS A LIMIT CONCEPT
%      pp. 74--78; divisions 28--29; notes 98--103
% ============================================================

\chapter{Totality as a Limit Concept}
\label{ch:6}

% [text to be added: pp. 74--75 --- division 28; notes 98--100]

% [text to be added: pp. 75--78 --- division 29; notes 101--103.
%  Last pages of the treatise.]


% ============================================================
%  APPENDIX: TABLE OF CATEGORIES (p. 79)
%  Announced in note 6.  Left to the end deliberately: its entries use
%  the vocabulary of the series-types fixed in chapters III--IV, and
%  should be rendered only once those renderings are settled.
% ============================================================

\chapter*{Table of Categories}
\addcontentsline{toc}{chapter}{Table of Categories}
\markboth{Table of Categories}{Table of Categories}

% [text to be added: p. 79 --- the KATEGORITAVLE, headed
%  "(Cf. Den menneskelige Tanke 1910)".  Four groups: I. Fundamental
%  Categories (1. Synthesis --- Relation; 2. Continuity --- Discontinuity;
%  3. Likeness --- Difference, with sub-items a--h, the last four braced
%  together as "Rationale Rækker"); II. Formal Categories (1. Identity;
%  2. Analogy; 3. Negation; 4. Rationality); III. Real Categories
%  (1. Causality; 2. Totality; 3. Development); IV. Ideal Categories
%  (Value-concepts).]

\end{document}
